% ============================================================
% f2012_guangdong_science_q17_histogram.tex
% 2012 广东卷（理科）第 17 题 频率分布直方图（TikZ 矢量重绘）
% 组距 10；各组 [40,50)…[90,100] 的“频率/组距”依次为
% 0.006, 0.006, 0.010, 0.054, x, 0.006，其中 x=0.018
% ============================================================
\begin{tikzpicture}[
  x=1cm, y=1cm, line width=0.6pt, font=\footnotesize, >=Stealth
]
  \path[draw=none] (-0.75,-0.75) rectangle (8.05,7.15);
  % 坐标轴
  \draw[->] (-0.35,0) -- (7.55,0) node[below] {成绩};
  \draw[->] (0,-0.05) -- (0,6.55);
  \node[right, inner sep=3pt] at (0.10,6.05) {$\dfrac{\text{频率}}{\text{组距}}$};
  % 断轴符号
  \draw (-0.30,0) -- (-0.20,0.15) -- (-0.10,0) -- (0.00,0.15);

  % 六个矩形，高度按 100 倍纵坐标绘制
  \draw (0.5,0) rectangle (1.5,0.6);
  \draw (1.5,0) rectangle (2.5,0.6);
  \draw (2.5,0) rectangle (3.5,1.0);
  \draw (3.5,0) rectangle (4.5,5.4);
  \draw (4.5,0) rectangle (5.5,1.8);
  \draw (5.5,0) rectangle (6.5,0.6);

  % 虚线参考线
  \draw[dashed, line width=0.4pt] (0,5.4) -- (3.5,5.4);
  \draw[dashed, line width=0.4pt] (0,1.8) -- (4.5,1.8);
  \draw[dashed, line width=0.4pt] (0,1.0) -- (2.5,1.0);
  \draw[dashed, line width=0.4pt] (0,0.6) -- (5.5,0.6);

  % 纵轴刻度值
  \node[left] at (0,5.4) {$0.054$};
  \node[left] at (0,1.8) {$x$};
  \node[left] at (0,1.0) {$0.010$};
  \node[left] at (0,0.6) {$0.006$};

  % 横轴刻度
  \draw (0,0) -- (0,-0.09);
  \node[below] at (0,-0.10) {$0$};
  \foreach \xx/\lab in {0.5/40,1.5/50,2.5/60,3.5/70,4.5/80,5.5/90,6.5/100}{
    \draw (\xx,0) -- (\xx,-0.09);
    \node[below] at (\xx,-0.10) {$\lab$};
  }

\end{tikzpicture}
